\scanpage{184}
% manuscript position: top
\mspage{100}
% manuscript position: 1–5
ce qui prouve $\mathrm{PF}_\pi1$. Prouvons $\mathrm{PF}_\pi2$, i.e. soient $x\xrightarrow{u}F$, $y\xrightarrow{v}F$ deux objets de $C_{/F}$, et $(\deleted{\unsure{u}}f_\alpha)_{\alpha\in A}$ une famille de cardinal $\leq\pi$ de flèches de $u$ dans $v$.
% manuscript position: left/4–5
\diagram{\begin{tikzcd}[ampersand replacement=\&,column sep=small,row sep=small]
x\arrow[rr,"f_\alpha"]\arrow[dr,"u"'] \& {} \& y\arrow[dl,"v"]\\
{} \& F \& {}
\end{tikzcd}}
% manuscript position: 5–7
On sait que $u,v$ proviennent de $\deleted{\unsure{f}}u_i:x\longrightarrow F_i$, $v_i:y\longrightarrow F_i$ pour $i$ assez grand.
% manuscript position: left/7–9
\diagram{\begin{tikzcd}[ampersand replacement=\&,column sep=small,row sep=small]
x\arrow[rr,"f_\alpha"]\arrow[dr,"u_i"'] \& {} \& y\arrow[dl,"v_i"]\\
{} \& F_i \& {}
\end{tikzcd}}
% manuscript position: 7–16
Pour $\unsure{\forall|\mathrm{tt}}\,\alpha$, on peut choisir $i_\alpha\geq i$ tel que $v_if_\alpha$ et $u_i$ deviennent égaux en composant avec $F_i\longrightarrow F_{i_\alpha}$, et comme $\operatorname{card} A\leq\pi$ et $I$ grd devant $\pi$, il \deleted{\unsure{faut}} \inserted{existe} un majorant $i'$ des $i_\alpha$. Au total, \deleted{\unsure{les}} \unsure{OPS|O.P.S.} (remplaçant $i$ par $i'$) que $v_i f_\alpha=u_i$ pour tout $\alpha$. Alors $u_i,v_i$ \deleted{\unsure{étant}} considérés comme objets de $C_{/F_i}$, les $\unsure{u_\alpha|f_\alpha}$ sont des flèches de $C_{/F_i}$ de $u_i$ dans $v_i$.
% manuscript position: right/16
\margin{$\unsure{=}$}

% manuscript position: 17–19
Comme $C_{/F_i}$ est grd devant $\pi$ (par l'hypothèse $F_i\in\operatorname{Ob}\operatorname{Ind}'_\pi(C)$), on trouve qu'il existe un égalisateur $(y,v_i)\xrightarrow{\unsure{g_i|g}}(z,w_i)$.
% manuscript position: left/20–22
\diagram{\begin{tikzcd}[ampersand replacement=\&,column sep=small,row sep=small]
x\arrow[r,shift left=1,"{(f_\alpha)}"]\arrow[r]\arrow[r,shift right=1]\arrow[dr,"u_i"'] \& y\arrow[r,"g"]\arrow[d,"v_i"] \& z\arrow[dl,"w_i"]\\
{} \& F_i \& {}
\end{tikzcd}}
% manuscript position: 20–27
Alors l'image de $g$ dans $C_{/F}$ (via $F_i\longrightarrow F$) égalise les flèches $f_\alpha$, en tant que flèches de $C_{/F}$, ce qui prouve $\mathrm{PF}_\pi2$ et achève la démonstration de 4.4.15, et par là aussi celle du th. 4.4.13, et du lemme 4.3. Par là est aussi achevée la dém. de la première assertion du lemme
% manuscript position: foot
% archive sheet number: 183
